\subsection{DDP\_cell\_porosity\_map.txt}
\index{Model!Hyperbolic!DDP}\label{sec_Model!Hyperbolic!DDP!cellPorosityMap}
\subsubsection{Purpose}
The Depth-Dependent Porosity (DDP) model requires that the variations in the storage porosity $\phi_\Omega$ with the water depth be specified on a cell-by-cell basis. This is the domain-based $\phi_\Omega(z)$ law in the original publication \cite{Guinot_al_2018}.
The code provides 5 types of pre-defined law:
\begin{itemize}
    \item Law type -2: the porosity law is specified as a sequence of couples $(z_i, \phi_i$), where the $z_i$ are ranked in ascending order. The $z_i$ are relative to the ground surface elevation, that is computed from the 2D mesh. The $z_i$ are the levels used in the tabulation of the law.
    \item Law type -1: the porosity law is specified as a sequence of couples $(z_i, \phi_i$), where the $z_i$ are ranked in ascending order. The $z_i$ are relative to the ground surface elevation, that is computed from the 2D mesh. The law is later tabulated using \textbf{ntabmax} intermediate levels between $z_1$ and $z_N$.
    \item Law type 0: the porosity law is specified by a sequence of elevations $(z_1, \ldots, z_N)$ in ascending order. The porosity is piecewise constant, equal to $i/N$ between $z_i$ and $z_i+1$ and to unity above $z_N$.  The $z_i$ are relative to the ground surface elevation as in Law type -1.
    \item Law type 1: the porosity is zero below an elevation $z_1$ relative to the ground elevation, and equal to $\phi_1$ above. 
    \item Law type 2: the porosity varies linearly from $\phi_1$ to $\phi_2$ between elevations $z_1$ and $z_2$ relative to the ground surface. It is zero below $z_1$ and $\phi_2$ above $z_2$.
    \item Law type 3: the porosity varies linearly from $\phi_1$ to $\phi_2$ between elevations $z_1$ and $z_2$ relative to the ground surface. It is zero below $z_1$ and $\phi_3$ above $z_2$.
\end{itemize}
If the laws are spatially distributed, each cell may obey a different law type.
%
\subsubsection{Law type specification format}
Five porosity law types are allowed in the current version of the software.\\

\begin{tabular}{p{1.5cm} | p{5.5cm} | p{8cm}} 
\hline
Law type & Parameters to be provided in file (separated using \texttt{<Tab>} characters) & Comments \\
\hline
\hline
-2 & -1 $N$ $\phi_1 \: z_1  \ldots \phi_N \: z_N$ &  $N$: number of $(z,\phi)$ couples. \\
\hline
-1 & -1 $N$ $\phi_1 \: z_1  \ldots \phi_N \: z_N$ &  $N$: number of $(z,\phi)$ couples. \\
\hline
0 & 0 $N$ $z_1 \ldots z_N$ & $N$ : number of points in the $\phi(z)$ law. $z_i$ : elevations relative to the ground level , by ascending order. \\
\hline
 1 & 1 $z_1$ $\phi_1$ & $z_1$: altitude (relative to local bottom elevation) above which $\phi=\phi_1$; $\phi = 0$ for $z < z_1$ \\
 \hline
 2 & 2 $z_1$ $z_2$ $\phi_1$ $\phi_2$ & Porosity $\phi$ is zero below $z_1$, varies linearly from $\phi_1$ to $\phi_2$ between $z_1$ and $z_2$, is equal to $\phi_2$ for $z > z_2$  \\
 \hline
 3 & 3 $z_1$ $z_2$ $\phi_1$ $\phi_2$ $\phi_3$ & Porosity $\phi$ is zero below $z_1$, varies linearly from $\phi_1$ to $\phi_2$ between $z_1$ and $z_2$, is equal to $\phi_3$ for $z > z_2$ \\
\hline
\end{tabular}
%
\subsubsection{Input file format}
%
File format for a uniform parameter set.  \texttt{\textbf{Bold}} : keyword, must appear as such in the file. \texttt{<Tab>}: tabulation character. \\
    \begin{tabular}{p{6cm} | p{9cm}}
    \hline
    Field(s) & Comment(s) \\
    \hline  \hline
    \texttt{\textbf{Unif} <Tab> \textbf{1}} & Flag is 1 because the parameter set is uniform \\ \hline
    \texttt{\textbf{NtabMax} <Tab> NtabMax} & NtabMax is the number of vertical discretization levels used for the $\phi_\Omega$(z) law \\ \hline
    \texttt{\textbf{==== Default Param}} & Section separator line, leave unchanged \\  \hline
    \texttt{\textbf{Uniform parameter set}} & Comment line, leave unchanged \\  \hline
    Parameter set     & The parameter set should obey the format given in the previous table \\ \hline
    \texttt{\textbf{==== Distrib}} & Section separator line, leave unchanged \\  \hline
    \end{tabular}
%
\vspace{10mm} \\
%
File format for a spatially distributed parameter set.  \texttt{\textbf{Bold}} : keyword, must appear as such in the file. \texttt{<Tab>}: tabulation character. \\
    \begin{tabular}{p{6cm} | p{9cm}}
    \hline
    Field(s) & Comment(s) \\
    \hline  \hline
    \texttt{\textbf{Unif} <Tab> \textbf{0}} & Flag is 0 because the parameter set is spatially variable \\ \hline
    \texttt{\textbf{NtabMax} <Tab> NtabMax} & NtabMax is the number of vertical discretization levels used for the $\phi_\Omega$(z) law \\ \hline
    \texttt{\textbf{==== Default Param}} & Section separator line, leave unchanged \\  \hline
    \texttt{\textbf{Uniform parameter set}} & Comment line, leave unchanged \\  \hline
    Parameter set     & The parameter set should obey the format given in the previous table \\ \hline
    \texttt{\textbf{==== Distrib}} & Section separator line, leave unchanged \\  \hline
    Parameter set [i]     & Parameter set, provided on a cell-by-cell basis. Index i runs sequentially from 1 to the number of cells. The parameter set should obey the format given in the Law type specification format table \\ \hline
    \end{tabular}
%
\subsubsection{Sample files}
%
\paragraph*{Example\#1: Uniform law type 2}
The porosity varies linearly from $\phi_1=0.$ to $\phi_2=1.$ between $z_1=3$ m and $z_2=10$ m.
\begin{verbatim}
Unif	1
NtabMax	50
==== Default Param
Type	Parameters(type)
2	3.	10.	0.	1.
==== Distrib
\end{verbatim}
%
\paragraph*{Example\#2: Distributed law type}
There are 3 cells in the model, the law types for cells 1, 2 and 3 are respectively 1, 0 and -1. The law type 0 in Cell 2 is given by 5 elevations. The Law type -1 in Cell 3 uses 2 $(z,\phi)$ couples.
\begin{verbatim}
Unif	0
NtabMax	50
==== Default Param
Type	Parameters(type)
2	3.	10.	0.	1.
==== Distrib
Type	Parameters(type)
1   2.3 0.5
0   5   0.1 0.4 0.7 1.2 2.5
-1  2   0.1 0.3 2.7 0.95
\end{verbatim}